
\medskip

%\emph{Les parties \textbf{A} et \textbf{B} sont indépendantes.\\
\emph{Les probabilités demandées seront données à $10^{-3}$ près.}

%Pour aider à la détection de certaines allergies, on peut procéder à un test sanguin dont le résultat est soit positif, soit négatif.
%
%Dans une population, ce test donne les résultats suivants :
%\begin{itemize}
%\item Si un individu est allergique, le test est positif dans 97\,\% des cas;
%\item Si un individu n'est pas allergique, le test est négatif dans 95,7\,\% des cas.
%\end{itemize}
%
%Par ailleurs, 20\,\% des individus de la population concernée présentent un test positif.
%
%On choisit au hasard un individu dans la population, et on note :
%
%\begin{itemize}
%\item[$\bullet~~$] $A$ l'évènement \og l'individu est allergique \fg{} ;
%\item[$\bullet~~$] $T$ l'évènement \og l'individu présente un test positif \fg.
%\end{itemize}
%
%On notera $\overline{A}$ et $\overline{T}$ les évènements contraires de $A$ et $T$.
%On appelle par ailleurs $x$ la probabilité de l'évènement $A$ : $x = p(A)$.

%\begin{minipage}{9cm}
\medskip

\textbf{Partie A}

\medskip

\begin{enumerate}
\item %Reproduire et compléter l'arbre ci-contre décrivant la situation, en indiquant sur chaque branche la probabilité correspondante.
D'après l'énoncé :

\begin{itemize}
\item[$\bullet~~$] $p_A(T) = 0,97$ ;
\item[$\bullet~~$] $p_{\overline{A}}\left(\overline{T}\right) = 0,957$ ;
\item[$\bullet~~$] $p(T) = 0,2$.
\end{itemize}

D'où l'arbre pondéré :

\begin{center}
\pstree[treemode=R,nodesepA=0pt,nodesepB=2.5pt,treesep = 1cm,levelsep=2.5cm]{\TR{}}
{\pstree{\TR{$A$~}\taput{$x$}}
	{\TR{$T$}\taput{0,97}
	\TR{$\overline{T}$}\tbput{0,03}
}
\pstree{\TR{$\overline{A}$~}\tbput{$1 - x$}}
{\TR{$T$}\taput{0,043}
\TR{$\overline{T}$}\tbput{0,957}
}
}
\end{center}

\item 
	\begin{enumerate}
		\item %Démontrer l'égalité : $p(T) = 0,927x + 0,043$.
D'après la loi des probabilités totales :
		
$p(T) = p(A \cap T) + p\left(\overline{A} \cap T\right) = p(A) \times p_{A}(T) + p\left(\overline{A}\right) \times p_{\overline{A}}(T) = x \times 0,97 + (1 - x) \times 0,043 = 0,97x + 0,043 -0,043x = 0,927x + 0,043$.
		\item %En déduire la probabilité que l'individu choisi soit allergique.
Comme $p(T) = 0,2 = 0,927x + 0,043 \iff 0,157 = 0,927x \iff \dfrac{0,157}{0,927} = x$.
		
Or $\dfrac{0,157}{0,927} \approx \np{0,1694}$ soit 0,169 au millième près.

$x  = p(A) \approx 0,169$.
	\end{enumerate}	
\item %Justifier par un calcul l'affirmation suivante :

%\og Si le test d'un individu choisi au hasard est positif, il y a plus de 80\,\% de chances que cet individu soit allergique \fg.
L'affirmation se traduit par : $p_T(A) > 0,8$.

Or $p_T(A) = \dfrac{p(T \cap A)}{p(T)} = \dfrac{p(A \cap T)}{p(T)} = \dfrac{p(A) \times p_A(T)}{p(T)} \approx \dfrac{0,169 \times 0,97}{0,2}$, soit 

$p_T(A) \approx \dfrac{\np{0,16393}}{0,2} \approx \np{0,81965}$, soit environ 81,97\,\% : l'affirmation est vraie.
\end{enumerate}
%\end{minipage}\hfill
%\begin{minipage}{5.5cm}

%\end{minipage}

\textbf{Partie B}

\medskip

%On réalise une enquête sur les allergies dans une ville en interrogeant $150$~habitants choisis au hasard, et on admet que ce choix se ramène à des tirages successifs indépendants avec remise.
%
%On sait que la probabilité qu'un habitant choisi au hasard dans cette ville soit allergique est égale à $0,08$.
%
%On note $X$ la variable aléatoire qui à un échantillon de $150$~habitants choisis au hasard associe le nombre de personnes allergiques dans cet échantillon.
%
%\medskip

\begin{enumerate}
\item %Quelle est la loi de probabilité suivie par la variable aléatoire $X$ ?

%Préciser ses paramètres.
Les tirages successifs étant indépendants et chaque personne ayant un probabilité d'être allergique égale à 0,08, $X$ suit une loi de Bernoulli de paramètres $n = 150$ et $p = 0,08$.
\item %Déterminer la probabilité que $20$~personnes exactement parmi les $150$ interrogées soient allergiques.
La calculatrice donne $p(X = 20) \approx \np{0,00820}$, soit 0,008 au millième près.
\item %Déterminer la probabilité qu'au moins 10\,\% des personnes parmi les $150$ interrogées soient allergiques.
La calculatrice donne $p(X \leqslant 14) \approx \np{0,7797}$, donc $p(X \geqslant 15) = 1 - p(X \leqslant 14) \approx 1 - \np{0,7797} \approx \np{0,2203}$ soit 0,220 au millième près.
\end{enumerate}

\bigskip

